\section{The network}
\label{sec:network}

\subsection{Siting rationale}
\label{sec:network-siting}

The \network{} was established to close a gap in the national observing system.
Coastal margins occupy a small fraction of land area and a disproportionate
share of the carbon exchange, and the existing network placed no tower within
15 kilometres of the shoreline \citep{sandvik2021coastalgap}. The stretch of
coast selected runs from the estuary at Kirkhavn in the north to the sand spit
at Roseholm in the south, and contains both of the land cover classes of
interest in interleaved patches of a size that a single tower can resolve.

Siting followed three constraints. Each tower needed a fetch of at least 300
metres of homogeneous cover in the prevailing wind sector, which for this coast
is 200 to 260 degrees. Each needed mains power or a solar array with 14 days of
winter autonomy. And the network as a whole needed at least three towers of each
land cover class so that a class mean would carry a usable standard error.

Table~\ref{tab:site-summary} lists the resulting sites.

\begin{table}[t]
\centering
\small
\caption{Sites of the \network. Elevation is metres above the ordnance datum.
Fetch is the shortest distance to a cover change within the prevailing sector.
The canopy height at the saltmarsh sites is the seasonal maximum reached in
August.}
\label{tab:site-summary}
\begin{tabular}{llrrrrl}
\toprule
Site & Cover & Elev. (m) & Tower (m) & Canopy (m) & Fetch (m) & Power \\
\midrule
MC01 & saltmarsh & 1.8  & 4.0 & 0.55 & 940  & mains \\
MC02 & saltmarsh & 2.1  & 4.0 & 0.61 & 720  & mains \\
MC03 & dune      & 8.4  & 6.0 & 0.90 & 480  & solar \\
MC04 & dune      & 11.2 & 6.0 & 1.15 & 610  & solar \\
MC05 & saltmarsh & 1.4  & 4.0 & 0.48 & 1150 & solar \\
MC06 & dune      & 6.9  & 6.0 & 0.82 & 340  & solar \\
MC07 & saltmarsh & 0.9  & 4.0 & 0.52 & 830  & solar \\
MC08 & dune      & 14.6 & 8.0 & 1.30 & 520  & mains \\
\bottomrule
\end{tabular}
\end{table}

MC06 has the shortest fetch in the network at 340 metres, which is above the
constraint but not comfortably. Its footprint climatology, computed by the
method of Section~\ref{sec:proc-footprint}, places 12 percent of the annual flux
source area outside the intended cover class under stable stratification. Fluxes
from MC06 under stable conditions are flagged and excluded from the class means
of Section~\ref{sec:results-budget}.

\subsection{Operating period and data return}
\label{sec:network-return}

All eight towers ran from 1 January to 31 December 2024. Two extended outages
affected the network: a storm on 14 to 17 February took the solar sites offline
when battery voltage fell below the cutoff, and a firmware fault in the
datalogger at MC02 and MC04 corrupted 11 days of high frequency records in
September before it was diagnosed. Neither loss is correlated with the flux
magnitude and both are treated as missing at random in the gap filling of
Section~\ref{sec:proc-gapfill}.
